\documentclass[a4paper,12pt,oneside]{maki-notes}

\renewcommand{\LectureNotesTitle}{高等数学讲义示例文档}
\renewcommand{\AuthorInfo}{Maki Notes Example}
\renewcommand{\CoverTitle}{高等数学讲义示例文档}
\renewcommand{\CoverAuthorBlock}{%
	Maki Notes 模板演示
	\\[0.8em]
	\Large
	内容覆盖: 定理, 例题, 真题, 表格, 图像
	\\
}
\renewcommand{\CoverFooterBlock}{%
	用于验证类文件与样式文件的真实排版效果
	\\[0.5em]
	适合作为二次开发与内容扩写的起点
	\\[2.5em]
	二〇二六年四月
}

\begin{document}

\frontmatter
\MakeTitlePage

\chapter*{前言}
本文档不是最小示例, 而是一份内容较为丰富的模板演示稿. 它主要用于验证模板在真实讲义场景中的可用性. 其中, 类文件入口为 \texttt{maki-notes.cls}, 版式定义集中在 \texttt{maki-notes.sty}. 文档刻意混合了定义、定理、证明、练习、真题、表格、二维函数图像与三维示意图, 目的是覆盖模板中的主要排版能力.

建议把本文件当作二次开发的起点: 一方面它足够短, 便于快速替换章节与题目; 另一方面它又足够完整, 可以直接检查封面、目录、页眉页脚、盒子环境、数学公式与图形系统是否工作正常.

\tableofcontents

\mainmatter

\chapter{极限与连续的基本结构}

\section{从定义出发理解极限}

极限思想的价值不在于记住若干结论, 而在于理解“无限逼近”如何转化为可验证的量化关系. 一旦把“靠近”写成精确的不等式条件, 很多直觉性判断就能被严格证明.

\begin{definition}[函数极限]
设函数 $f(x)$ 在点 $x_0$ 的某去心邻域内有定义. 若存在常数 $A$, 使得对于任意给定的 $\varepsilon > 0$, 总存在 $\delta > 0$, 当
\[
0 < |x - x_0| < \delta
\]
时, 恒有
\[
|f(x) - A| < \varepsilon,
\]
则称 $A$ 为函数 $f(x)$ 在 $x \to x_0$ 时的极限, 记作
\[
\lim_{x \to x_0} f(x) = A.
\]
\end{definition}

\begin{theorem}[极限的唯一性]
若 $\lim_{x \to x_0} f(x)$ 存在, 则它唯一.
\end{theorem}

\begin{proof}
假设 $\lim_{x \to x_0} f(x) = A$ 与 $\lim_{x \to x_0} f(x) = B$, 且 $A \neq B$. 取
\[
\varepsilon = \dfrac{|A-B|}{2}.
\]
由极限定义, 分别存在 $\delta_1 > 0$ 与 $\delta_2 > 0$, 使得当 $0 < |x-x_0| < \delta_1$ 时有 $|f(x)-A| < \varepsilon$, 当 $0 < |x-x_0| < \delta_2$ 时有 $|f(x)-B| < \varepsilon$. 取 $\delta = \min\{\delta_1,\delta_2\}$, 则当 $0<|x-x_0|<\delta$ 时,
\[
|A-B| \le |A-f(x)| + |f(x)-B| < 2\varepsilon = |A-B|,
\]
矛盾. 故 $A=B$.
\end{proof}

\begin{note}
唯一性证明的结构非常典型: 先假设存在两个不同候选值, 再通过三角不等式把两边同时逼近的结论拼起来, 最后制造严格小于自身的矛盾. 这种证明方式在连续性、导数与积分的理论中都会反复出现.
\end{note}

\begin{example}
用定义证明
\[
\lim_{x \to 1} (2x^2 - x + 3) = 4.
\]
\end{example}

\textbf{解.} 先化简:
\[
|2x^2 - x + 3 - 4| = |2x^2 - x - 1| = |x-1||2x+1|.
\]
若再附加 $|x-1| < 1$, 则 $x \in (0,2)$, 从而 $|2x+1| < 5$. 因此只要令
\[
\delta = \min\left\{1,\dfrac{\varepsilon}{5}\right\},
\]
当 $0<|x-1|<\delta$ 时就有
\[
|2x^2 - x + 3 - 4| < 5|x-1| < 5\delta \le \varepsilon.
\]
故结论成立. \hfill $\square$

\section{连续性的语言}

\begin{definition}[连续]
若函数 $f(x)$ 在点 $x_0$ 处满足
\[
\lim_{x \to x_0} f(x) = f(x_0),
\]
则称 $f(x)$ 在 $x_0$ 处连续.
\end{definition}

\begin{property}[连续函数的四则运算]
若 $f,g$ 在 $x_0$ 处连续, 则 $f \pm g$, $fg$ 在 $x_0$ 处连续; 若进一步 $g(x_0)\neq 0$, 则 $\dfrac{f}{g}$ 在 $x_0$ 处连续.
\end{property}

\begin{exercise}
设 $f(x)=\dfrac{x^2-1}{x-1}$, 在 $x\neq 1$ 时有定义. 若要把它补成在 $x=1$ 连续的函数, 则应定义 $f(1) = $ \fillin[2.5cm].
\end{exercise}

\begin{exercise}
下列叙述中正确的是\selectbracket
\begin{enumerate}
\item[A.] 可导 $\Longrightarrow$ 连续
\item[B.] 连续 $\Longrightarrow$ 可导
\item[C.] 极限存在 $\Longrightarrow$ 连续
\item[D.] 有界 $\Longrightarrow$ 连续
\end{enumerate}
\end{exercise}

\chapter{导数与最值问题}

\section{导数的几何直观}

导数可以理解为“局部线性近似的斜率”. 在研究函数的增减性与最值时, 导数提供了连接图像直觉与代数判断的桥梁.

\begin{definition}[导数]
设 $f(x)$ 在点 $x_0$ 的某邻域内有定义. 若极限
\[
f'(x_0)=\lim_{\Delta x \to 0}\dfrac{f(x_0+\Delta x)-f(x_0)}{\Delta x}
\]
存在, 则称 $f$ 在 $x_0$ 处可导, 并称该极限为 $f$ 在 $x_0$ 处的导数.
\end{definition}

\begin{theorem}[Fermat 引理]
若 $f$ 在 $x_0$ 处可导且在该点取得极值, 则 $f'(x_0)=0$.
\end{theorem}

\begin{corollary}
可导函数若在开区间内取得极大值或极小值, 则对应驻点必须满足导数为零. 但导数为零并不必然意味着极值存在.
\end{corollary}

\begin{longexample}
求函数
\[
f(x)=x^3-3x^2+2
\]
在区间 $[0,3]$ 上的单调区间与最值.

\tcblower

\textbf{解.} 先求导:
\[
f'(x)=3x^2-6x=3x(x-2).
\]
于是当 $x\in(0,2)$ 时, $f'(x)<0$; 当 $x\in(2,3)$ 时, $f'(x)>0$. 因而 $f$ 在 $(0,2)$ 上单调递减, 在 $(2,3)$ 上单调递增.

接着比较端点与驻点的函数值:
\[
f(0)=2,\qquad f(2)=-2,\qquad f(3)=2.
\]
因此区间最小值为 $-2$, 在 $x=2$ 处取得; 区间最大值为 $2$, 在 $x=0$ 与 $x=3$ 处取得.
\end{longexample}

\begin{realexam}[2024 模拟卷]
设函数 $f(x)$ 在 $[0,1]$ 上连续, 在 $(0,1)$ 内可导, 且 $f(0)=0$, $f(1)=2$. 证明: 存在 $\xi \in (0,1)$, 使得
\[
f'(\xi)=2.
\]
\end{realexam}

\textbf{证.} 这正是 Lagrange 中值定理的直接应用. 因为
\[
\dfrac{f(1)-f(0)}{1-0} = \dfrac{2-0}{1}=2,
\]
所以存在 $\xi \in (0,1)$ 使得 $f'(\xi)=2$. \hfill $\square$

\section{学习进度的表格化展示}

仅有公式并不足以验证模板的实用性, 表格排版同样重要. 下表用来展示讲义编写计划与每章侧重点.

\begin{center}
\renewcommand{\arraystretch}{1.35}
\begin{tabular}{ccccc}
\toprule
\diagbox{周次}{模块} & 极限与连续 & 导数与微分 & 积分学 & 空间解析几何 \\
\midrule
\multirow{2}{*}{第 1 周} & 定义与性质 & 导数概念 & 原函数 & 向量代数 \\
 & 典型证明 & 几何意义 & 基本公式 & 点线面关系 \\
\midrule
\multirow{2}{*}{第 2 周} & 连续函数 & 中值定理 & 定积分方法 & 曲面方程 \\
 & 极限计算 & 单调与最值 & 应用问题 & 参数方程 \\
\midrule
第 3 周 & 综合训练 & 真题分析 & 综合训练 & 图形直观 \\
\bottomrule
\end{tabular}
\captionof{table}{三周讲义组织示例}
\end{center}

\chapter{积分与建模表达}

\section{从面积到累计量}

\begin{axiom}[Newton--Leibniz 公式]
设 $f$ 在 $[a,b]$ 上连续, 且 $F'(x)=f(x)$, 则
\[
\int_a^b f(x)\,\mathrm{d}x = F(b)-F(a).
\]
\end{axiom}

\begin{example}
计算
\[
I=\int_0^1 (3x^2+2x+1)\,\mathrm{d}x.
\]
\end{example}

\textbf{解.} 原函数可以取
\[
F(x)=x^3+x^2+x.
\]
于是
\[
I=F(1)-F(0)=3.
\]

\begin{proposition}
若 $f(x)\ge 0$ 在区间 $[a,b]$ 上恒成立, 则
\[
\int_a^b f(x)\,\mathrm{d}x \ge 0.
\]
\end{proposition}

\begin{note}
在应用问题中, 定积分常常表示“累计量”, 例如位移、质量、总成本或总收益. 讲义若要写得自然, 可以先说明被积函数代表什么局部变化, 再说明积分为何表示总体累计.
\end{note}

\begin{longrealexam}[2023 真题改编]
设函数 $f(x)$ 在 $[0,1]$ 上连续, 且满足
\[
f(x)=1+\int_0^x (x-t)f(t)\,\mathrm{d}t.
\]
求 $f''(x)$ 与 $f(x)$ 的关系, 并由此解出 $f(x)$.

\tcblower

\textbf{解.} 对原式求导:
\[
f'(x)=\int_0^x f(t)\,\mathrm{d}t.
\]
再对上式求导, 得
\[
f''(x)=f(x).
\]
由原积分方程可知 $f(0)=1$, 又由 $f'(x)=\int_0^x f(t)\,\mathrm{d}t$ 得 $f'(0)=0$. 因此所求函数满足常系数微分方程
\[
f''-f=0,\qquad f(0)=1,\qquad f'(0)=0.
\]
解得
\[
f(x)=\dfrac{e^x+e^{-x}}{2}=\cosh x.
\]
\end{longrealexam}

\chapter{图像与空间直观}

\section{函数图像}

下面给出一个二维函数图像, 用来检查 `pgfplots` 配置、颜色、坐标轴标签位置与图例是否正常.

\begin{center}
\begin{tikzpicture}
\begin{axis}[
	width=11cm,
	height=6.5cm,
	axis lines=middle,
	xlabel={$x$},
	ylabel={$y$},
	samples=120,
	domain=-2:4,
	grid=major,
	grid style={dashed, gray!30},
	legend pos=north east,
	ymin=-1.5,
	ymax=4.5
]
\addplot[thick, color=LogicColor]{x^2*exp(-x/2)};
\addlegendentry{$y=x^2e^{-x/2}$}
\addplot[thick, dashed, color=ExColor]{exp(-x/2)*(2*x - x^2/2)};
\addlegendentry{$y'=e^{-x/2}(2x-\dfrac{x^2}{2})$}
\end{axis}
\end{tikzpicture}
\captionof{figure}{函数及其导函数示意}
\end{center}

\section{空间向量示意}

三维图形常用来表达法向量、投影与曲面方向关系. 下面的示意图用于验证 `tikz-3dplot` 与颜色系统协同工作正常.

\begin{center}
\begin{tikzpicture}[tdplot_main_coords, scale=1.25]
\draw[-{Stealth}, thick] (0,0,0) -- (3.4,0,0) node[anchor=north east]{$x$};
\draw[-{Stealth}, thick] (0,0,0) -- (0,3.4,0) node[anchor=north west]{$y$};
\draw[-{Stealth}, thick] (0,0,0) -- (0,0,3.6) node[anchor=south]{$z$};

\draw[-{Stealth}, very thick, color=LogicColor] (0,0,0) -- (2.2,1.0,2.1) node[midway, above left]{$\bm a$};
\draw[-{Stealth}, very thick, color=ExColor] (0,0,0) -- (1.0,2.6,0.8) node[midway, right]{$\bm b$};
\draw[-{Stealth}, very thick, dashed, color=NoteColor] (0,0,0) -- (-1.5,1.1,2.2) node[anchor=south west]{$\bm a \times \bm b$};

\draw[pattern=north east lines, pattern color=gray!50, opacity=0.25]
	(0,0,0) -- (2.2,1.0,2.1) -- (3.2,3.6,2.9) -- (1.0,2.6,0.8) -- cycle;

\node[anchor=north east] at (0,0,0) {$O$};
\end{tikzpicture}
\captionof{figure}{向量张成平行四边形的空间示意}
\end{center}

\section{综合练习}

\begin{exercise}
设 $z=x^2y+e^{xy}$, 则
\[
\dfrac{\partial^2 z}{\partial x \partial y} = \fillin[4cm].
\]
\end{exercise}

\begin{exercise}
若平面 $\Pi$ 过点 $(1,-2,3)$ 且法向量为 $\bm n=(2,1,-1)$, 则平面方程可写为 \fillin[5cm].
\end{exercise}

\begin{note}
到这里为止, 模板中的标题页、目录、页眉页脚、盒子环境、跨页例题、题型辅助命令、表格与图形模块都已经被真实内容覆盖了一遍. 如果你要继续扩展模板, 最稳妥的方式就是先在本文件中加一个新章节或新环境, 确认编译稳定后再向正式讲义迁移.
\end{note}

\end{document}
